\documentclass[11pt]{article}
\usepackage[margin=1in]{geometry}
\usepackage{amsmath,amssymb,amsthm,hyperref}
\newtheorem{proposition}{Proposition}
\title{Growth Regimes with Endogenous Task Automation}
\author{Exploratory note for Oliver Pardo}
\date{22 November 2026}
\begin{document}
\maketitle
I retain the task technology and equilibrium conventions of
\texttt{model\_note.tex}. This note establishes exact balanced-growth equilibria
under competition and conditional limiting restrictions under monopoly.
It does not establish an exhaustive classification of dynamic equilibria.

\section{Normalization and the two competitive thresholds}
Let $H=AN$, $\dot H/H=n+\gamma$, $k=K/H$, $x=X/H$, and
$z(x)=\mathcal Z(1,x)$. Define
\[
 R_* =\rho+\gamma+\delta,\qquad
 a=\left(\frac{\alpha}{R_*}\right)^{1/(1-\alpha)},\qquad
 B_a=\frac1{(1-\alpha)a^\alpha}.
\]
I assume $\rho>n$, $\gamma>0$, and the original interior parameter ranges.
Here $B_a$ is an efficiency threshold, not an upper bound. At a balanced-growth
allocation with finite $k,x$, the household Euler equation requires
$r=\rho+\gamma$. The capital condition then gives $k=a z(x)$.
Under competitive AI supply, $B$ stays at $B_0$ and $p_X=1/B_0$.

Let $m$ be the endogenous task cutoff, and define
\[
 D_\sigma(m)=\int_0^m e^{-\beta(\sigma-1)i}\,di,\quad
 S(m)=1-m+D_\sigma(m)e^{\beta(\sigma-1)m},\quad
 c_\sigma=\left[D_\sigma(1)\right]^{1/(\sigma-1)}.
\]
At $\sigma=1$, use $c_1=e^{-\beta/2}$.
The marginal product $z_x$ along the optimal allocation decreases strictly
from $1$ to $c_\sigma$. Indeed,
\[
 \frac{d\log z_x}{dm}=-\frac{\beta(1-m)}{S(m)}<0.
\]
This formula also holds at unit elasticity, where $S=1$.
Define the second threshold
\begin{equation}
 B_g^{\mathrm{mc}}(\sigma)=\frac{B_a}{c_\sigma}>B_a.
 \label{eq:regimes-comp-threshold}
\end{equation}

\begin{proposition}[Exact competitive balanced-growth allocations]
If $B_0\leq B_a$, there is a competitive balanced-growth equilibrium with
$k=a$, $X=0$ and no research. If
$B_a<B_0<B_g^{\mathrm{mc}}(\sigma)$, there is exactly one competitive
balanced-growth allocation with finite positive $k,x$. Its cutoff solves
\begin{equation}
 z_x(m)=\frac{B_a}{B_0},\qquad
 x=\frac{D_\sigma(m)e^{\sigma\beta m}}{1-m},\qquad k=a z(x).
 \label{eq:regimes-coexistence}
\end{equation}
Starting at the corresponding $K_0=kH_0$, these allocations extend to
infinite-horizon equilibria. Output per worker and wages grow at $\gamma$.
These statements do not claim convergence from arbitrary initial capital.
\end{proposition}
\begin{proof}
When $X=0$, $z=1$ and the capital condition gives $k=a$.
The right marginal value of AI services is $(1-\alpha)a^\alpha=1/B_a$,
so zero adoption maximizes final-firm profit exactly when $B_0\leq B_a$.
With positive adoption, substitute $k/z=a$ into
$(1-\alpha)(k/z)^\alpha z_x=1/B_0$. Strict monotonicity and the endpoint
values of $z_x$ prove existence and uniqueness of the positive finite solution.

For either allocation, let $y=k^\alpha z^{1-\alpha}$ and $u=x/B_0$.
Set $C/H=y-u-(\delta+n+\gamma)k$. The AI revenue share
$s_X=xz_x/z$ belongs to $[0,1)$, and
\[
 \frac{C}{K}=\frac{R_*}{\alpha}
 [1-(1-\alpha)s_X]-\delta-n-\gamma>\rho-n>0.
\]
The weak bound suffices in the zero-adoption case as well.
The resource constraint and Euler equation hold. The household transversality
term is proportional to $e^{-(\rho-n)t}$ and converges to zero.
Concavity of production and log utility makes the household first-order
conditions sufficient under the usual no-Ponzi restriction. Competitive
AI providers earn zero profit. Under the same nonappropriable-research
benchmark as the original paper, private research is zero.
\end{proof}

The constant labor share in the positive-adoption equilibrium is
$s_L=(1-\alpha)z_H/z>0$. The no-adoption share is $1-\alpha$.
There is no positive finite balanced-growth solution at
$B_0\geq B_g^{\mathrm{mc}}$. At equality, the finite-cutoff solution escapes
to $m=1$. The critical dynamic case requires separate analysis.

\section{Conditional AI-dominated growth under either market structure}
Consider a path with $X/H\to\infty$ and $B\to B_\infty>0$.
The task result gives $m\to1$, $Z/X\to c_\sigma$ and $Z_X\to c_\sigma$.
For competitive pricing use $\lambda=1$. For a monopoly that globally solves
its static supply problem at each date, the limiting inverse-demand elasticity
is $\alpha$, so use $\lambda=1-\alpha$. The limiting supply condition is
\[
 \frac1{B_\infty}=\lambda(1-\alpha)(K/Z)^\alpha c_\sigma.
\]
Consequently,
\begin{equation}
 r\longrightarrow\mathcal R^j(B_\infty)
 =\alpha[B_\infty\lambda_j(1-\alpha)c_\sigma]^{(1-\alpha)/\alpha}-\delta,
 \quad \lambda_{\mathrm{mc}}=1,\quad\lambda_{\mathrm{mon}}=1-\alpha.
 \label{eq:regimes-returns}
\end{equation}
If $C/K$ converges to a positive constant, the Euler equation implies
\[
 g_{Y/N}=g_{K/N}=\mathcal R^j(B_\infty)-\rho.
\]
This refers to limiting average growth, and to instantaneous growth only
when those rates converge. Permanently faster growth requires
\begin{equation}
 B_\infty>B_g^j(\sigma),\qquad
 B_g^{\mathrm{mon}}(\sigma)=\frac{B_a}{(1-\alpha)c_\sigma}
 =\frac{B_g^{\mathrm{mc}}(\sigma)}{1-\alpha}.
 \label{eq:regimes-mon-threshold}
\end{equation}
These are conditional threshold identities, not dynamic existence theorems.
In particular, $B_\infty$ cannot yet be replaced by $\overline B$ without
establishing attainment along the relevant monopoly trajectory.

The wage identity $w/A=p_Xe^{\beta m}$ gives limiting average wage growth
$\gamma$, since the AI price converges to a positive constant.
If output per worker grows faster than $\gamma$, labor's income share
vanishes. AI-industry revenue converges to share $1-\alpha$ of gross output.
Competitive operating profit is zero. Under monopoly, compute expenditure
and operating profit converge to shares $(1-\alpha)^2$ and
$\alpha(1-\alpha)$ of gross output, respectively. These operating-profit
limits do not identify net profit without accounting for research expenditure.

\section{What changes with the elasticity}
The coefficient $c_\sigma$ is a power mean of $e^{-\beta i}$, with order
$\sigma-1$. It therefore increases strictly with $\sigma$, with limits
\[
 \lim_{\sigma\downarrow0}c_\sigma=\frac{\beta}{e^\beta-1},\qquad
 \lim_{\sigma\to\infty}c_\sigma=1.
\]
For completeness, writing $V=-\beta i$ and
$L(q)=\log\int_0^1e^{qV}\,di$, convexity gives
$qL'(q)-L(q)\geq0$ because $L(0)=0$. Strict convexity for nonconstant $V$
gives strict increase of $L(q)/q=\log c_\sigma$, including its continuous
extension at $q=0$.
Thus the conditional return in~\eqref{eq:regimes-returns} rises with
$\sigma$, and both growth thresholds fall. Complementarity between tasks
does not rule out AI-dominated growth when every task is automatable.
Wages nevertheless have limiting average growth $\gamma$ in all the
regular AI-dominated cases described above. The finite wage-growth maximizer
in the original fixed-weight model is not reproduced by this benchmark.

\section{Monopoly cases that remain to be resolved}
There is one exact zero-research construction. If $B_0\leq\overline B\leq B_a$,
set $k=a$, $X=M=0$, and $B=B_0$. Facing these aggregate paths, the developer
cannot earn positive operating profit at any attainable efficiency, since
$P(0)=1/B_a\leq1/B$. Any positive research expenditure lowers profit.
The competitive no-adoption construction therefore also supplies a monopoly
equilibrium. If $\overline B>B_a$, this argument no longer works. Research
that crosses the adoption threshold must be evaluated globally before declaring
zero research optimal.

For a finite positive limiting allocation with $r\to\rho+\gamma$, the
static monopoly first-order condition instead requires
\begin{equation}
 z_x(m)[1-e_X(m)]=\frac{B_a}{B_\infty}.
 \label{eq:regimes-mon-finite}
\end{equation}
The left side starts at $1$ and converges to $(1-\alpha)c_\sigma$.
These endpoints alone do not establish monotonicity or uniqueness.
Every root must be checked against all alternative quantities at its own
$k=a z(x)$. The developer's dynamic optimum, the limit of $B$, research costs
and household feasibility must then be verified. I do not label these roots
equilibria before those checks.

The unresolved tasks are global static monopoly selection, dynamic research
optimality, existence of AI-dominated paths, and the critical threshold cases.
No transition simulations have been run for this note.
\end{document}
